\section{Related Work}

\subsection{Theoretical Foundations of SVC}

SVC traces its theoretical roots to Vapnik's SVM framework~\cite{scholkopf2002} and its extension to density estimation. Tax and Duin~\cite{tax1999,tax2004} proposed SVDD, providing the direct optimization form for SVC: seeking the minimal hypersphere that tightly envelops the data in feature space. Ben-Hur \etal~\cite{benhur2001} extended SVDD into a complete clustering framework by introducing a connectivity-checking mechanism for multi-cluster label assignment.

Solving the dual QP of Eq.~\eqref{eq:svc_qp}, data points are classified into three categories by their $\alpha_i$ values: (i)~interior points ($\alpha_i\approx0$), inside the sphere and not participating in its definition; (ii)~support vectors (SV, $0<\alpha_i<C$), on the sphere surface and defining its geometry; (iii)~bounded support vectors (BSV, $\alpha_i\approx C$), outside the sphere, representing outliers or noise.

For any point $y$, its squared feature-space distance to the sphere center $a$ is:
\begin{equation}
\begin{split}
R^2(y)=\|\phi(y)-a\|^2=K(y,y) \\
-2\sum_i\alpha_i K(y,x_i)+\sum_{i,j}\alpha_i\alpha_j K(x_i,x_j)
\end{split}
\end{equation}

The sphere radius $R$ is taken as the median $R^2$ value over support vectors. Connectivity labeling operates by checking, for each pair of spatially close points $(x_i,x_j)$, whether the sampled points $y_t=x_i+t(x_j-x_i)$ along the connecting segment satisfy $R^2(y_t)\leq R^2$ for all $t\in[0,1]$.

\subsection{Improvements and Variants of SVC}

Since SVC's introduction, research has advanced along three main directions.

\textbf{Connectivity labeling efficiency.} The original SVC requires segment sampling for all $O(N^2)$ point pairs, yielding $O(N^3)$ complexity. Lee and Lee~\cite{lee2005} proposed an adjacency-matrix and graph-connected-component method reducing it to $O(N^2)$. Yang \etal~\cite{yang2002} introduced proximity graph modeling, limiting the number of point pairs needing checking. Our SVC implementation adopts $k$-NN graph filtering, checking only each point's $k$ spatial nearest neighbors, reducing to $O(kN)$.

\textbf{Dynamic characterization of cluster labels.} Lee and Lee~\cite{lee2007} proposed tracking the stability of cluster structures across different $(q,C)$ values to assist parameter selection. The core observation is that cluster structure remains relatively stable within appropriate parameter intervals. However, this method still requires human interpretation of ``stability intervals'' and has not achieved full automation.

\textbf{Large-scale extensions.} Researchers have explored acceleration strategies including core-set approximation~\cite{tsang2006}, random feature mapping~\cite{rahimi2008}, and convex clustering equivalence~\cite{hofmann2008}, focusing primarily on computational efficiency rather than parameter automation.

\subsection{Parameter Selection in Kernel Methods}

Parameter selection for kernel methods is a broad research problem. In supervised SVMs, parameters (e.g., $C$ and $\gamma$) are typically selected via $k$-fold cross-validation using labels~\cite{scholkopf2002}. In unsupervised settings, the problem is far more difficult.

Scholkopf~\etal~\cite{scholkopf2001} introduced the $\nu$-parameter in one-class SVM to control the outlier fraction, providing a more intuitive alternative to the raw $C$ parameter. $\nu$ is defined as an upper bound on the outlier fraction and a lower bound on the SV fraction. ASVC inherits this parameterization, replacing $C$ with $C=1/(\nu N)$, giving the parameter a probabilistic interpretation.

Wang and Zhu~\cite{wang2008} proposed a stability-based parameter selection method for SVC, identifying effective parameters by tracking the stability intervals of cluster counts as $q$ varies. The underlying assumption is that clustering within a stability interval has physical meaning. However, this assumption fails on certain datasets---as shown in our experiments, OverlappingGaussians yields no effective clustering for manual SVC at any $q$, hence no ``stability interval'' exists.

\subsection{Clustering Evaluation Metrics}

Clustering evaluation falls into external and internal categories. External metrics require labels: Adjusted Rand Index (ARI)~\cite{hubert1985}, Normalized Mutual Information (NMI)~\cite{strehl2002}, Homogeneity, and Completeness~\cite{rosenberg2007}. We use ARI and NMI as primary quantitative metrics.

Internal metrics (Silhouette~\cite{rousseeuw1987}, Davies-Bouldin~\cite{davies1979}, Calinski-Harabasz~\cite{calinski1974}) do not require labels but are based on Euclidean distance measures that do not fully match SVC's feature-space spherical geometry model. Therefore, we do not use them to guide parameter selection, instead designing an SVC-specific scoring function.
